%==================================================================
% Ini adalah bab 4
% Silahkan edit sesuai kebutuhan, baik menambah atau mengurangi \section, \subsection
%==================================================================

\chapter[HASIL DAN PEMBAHASAN]{\\ HASIL DAN PEMBAHASAN}

\section{Deskripsi Hasil Penelitian}
\begin{sectioncontent}
    \hspace{\parindent}Penelitian ini dilakukan untuk membandingkan performa YOLOv8 Nano dan YOLO11 Nano pada tugas klasifikasi penyakit daun tomat. Kedua model dilatih menggunakan dataset yang sama dan konfigurasi pelatihan yang dibuat seragam agar perbandingan dapat dilakukan secara adil. Setelah proses pelatihan selesai, model terbaik dari masing-masing arsitektur diekspor ke format ONNX untuk diuji pada perangkat \textit{edge} menggunakan Raspberry Pi 5.

    Pengujian dilakukan dengan dua fokus utama, yaitu akurasi klasifikasi dan efisiensi inferensi. Akurasi klasifikasi digunakan untuk menilai kemampuan model dalam memprediksi kelas penyakit daun tomat secara benar, sedangkan efisiensi inferensi digunakan untuk menilai kecepatan model ketika dijalankan pada perangkat \textit{edge}. Dengan demikian, model yang dianggap lebih baik tidak hanya dilihat dari nilai akurasi, tetapi juga dari kemampuannya untuk berjalan secara efisien pada perangkat dengan sumber daya terbatas.
\end{sectioncontent}

\section{Dataset Pengujian}
\begin{sectioncontent}
    \hspace{\parindent}Dataset pengujian yang digunakan terdiri dari 720 citra daun tomat yang terbagi ke dalam enam kelas. Setiap kelas memiliki jumlah data yang sama, yaitu 120 citra. Data uji ini tidak digunakan dalam proses pelatihan model, sehingga dapat digunakan untuk mengevaluasi performa akhir kedua model secara lebih objektif. Pembagian jumlah data pada masing-masing kelas dapat dilihat pada Tabel \ref{tab:jumlah-data-uji}.

    \begin{table}[H]
        \centering
        \caption{Jumlah Data Uji pada Setiap Kelas}
        \label{tab:jumlah-data-uji}
        \begin{tabular}{|l|c|}
            \hline
            \textbf{Kelas} & \textbf{Jumlah Data Uji} \\
            \hline
            Tomato Early Blight Leaf & 120 \\
            \hline
            Tomato Healthy Leaf & 120 \\
            \hline
            Tomato Leaf Yellow Curl Virus & 120 \\
            \hline
            Tomato Mold Leaf & 120 \\
            \hline
            Tomato Septoria Leaf Spot & 120 \\
            \hline
            Tomato Leaf Late Blight & 120 \\
            \hline
            \textbf{Total} & \textbf{720} \\
            \hline
        \end{tabular}
    \end{table}

    Berdasarkan Tabel \ref{tab:jumlah-data-uji}, dataset pengujian yang digunakan bersifat seimbang karena setiap kelas memiliki jumlah citra yang sama. Kondisi ini membantu proses interpretasi hasil evaluasi karena nilai akurasi tidak terlalu dipengaruhi oleh dominasi kelas tertentu. Meskipun demikian, penelitian ini tetap menggunakan metrik \textit{precision}, \textit{recall}, dan \textit{F1-score} untuk melihat performa model secara lebih rinci pada setiap kelas.
\end{sectioncontent}

\section{Hasil Pelatihan Model}
\begin{sectioncontent}
    \hspace{\parindent}Pelatihan model dilakukan menggunakan YOLOv8 Nano dan YOLO11 Nano versi klasifikasi. Kedua model dilatih menggunakan dataset yang sama dengan ukuran citra masukan sebesar 224 x 224 piksel. Untuk menjaga keadilan eksperimen, parameter utama seperti jumlah \textit{epoch}, ukuran citra, dan pembagian dataset dibuat sama pada kedua model.

    \begin{table}[H]
        \centering
        \caption{Konfigurasi Pelatihan Model}
        \label{tab:konfigurasi-pelatihan-bab4}
        \begin{tabular}{|l|l|}
            \hline
            \textbf{Parameter} & \textbf{Nilai} \\
            \hline
            Model 1 & YOLOv8 Nano Classification \\
            \hline
            Model 2 & YOLO11 Nano Classification \\
            \hline
            Ukuran citra & 224 x 224 piksel \\
            \hline
            Jumlah \textit{epoch} & 50 \\
            \hline
            Batch size & 16 \\
            \hline
            Perangkat pelatihan & NVIDIA RTX 4060 \\
            \hline
            Perangkat pengujian \textit{edge} & Raspberry Pi 5 RAM 8 GB \\
            \hline
        \end{tabular}
    \end{table}

    Hasil pelatihan dari masing-masing model menghasilkan bobot terbaik yang kemudian digunakan untuk proses validasi dan pengujian. Visualisasi hasil pelatihan, seperti grafik perubahan \textit{loss} dan akurasi selama proses pelatihan, dapat dilihat pada Gambar \ref{fig:training-yolov8} dan Gambar \ref{fig:training-yolo11}.

    % Siapkan gambar dari folder runs/classify/train*/results.png
    \begin{figure}[H]
        \centering
        \includegraphics[width=0.9\textwidth]{results-yolov8n.png}
        \caption{Grafik hasil pelatihan YOLOv8 Nano}
        \label{fig:training-yolov8}
    \end{figure}

    \begin{figure}[H]
        \centering
        \includegraphics[width=0.9\textwidth]{results-yolo11n.png}
        \caption{Grafik hasil pelatihan YOLO11 Nano}
        \label{fig:training-yolo11}
    \end{figure}

    Gambar hasil pelatihan digunakan untuk melihat perkembangan performa model selama proses pelatihan. Apabila nilai \textit{loss} menurun dan metrik evaluasi meningkat secara stabil, maka proses pelatihan dapat dianggap berjalan dengan baik. Namun, apabila terdapat perbedaan besar antara performa data latih dan data validasi, maka perlu dilakukan analisis lebih lanjut terhadap kemungkinan \textit{overfitting}.
\end{sectioncontent}

\section{Hasil Pengujian Klasifikasi}
\begin{sectioncontent}
    \hspace{\parindent}Pengujian klasifikasi dilakukan menggunakan data uji yang sama untuk kedua model. Hasil prediksi model dibandingkan dengan label sebenarnya untuk memperoleh nilai akurasi, \textit{precision}, \textit{recall}, dan \textit{F1-score}. Pada tahap pengujian, penamaan kelas pada data uji disesuaikan dengan penamaan kelas yang digunakan pada proses pelatihan agar hasil evaluasi tidak dipengaruhi oleh perbedaan nama folder kelas. Ringkasan hasil pengujian klasifikasi dapat dilihat pada Tabel \ref{tab:hasil-klasifikasi}.

    \begin{table}[H]
        \centering
        \caption{Perbandingan Hasil Pengujian Klasifikasi}
        \label{tab:hasil-klasifikasi}
        \renewcommand{\arraystretch}{1.2}
        \begin{tabularx}{\textwidth}{|l|>{\centering\arraybackslash}X|>{\centering\arraybackslash}X|>{\centering\arraybackslash}X|>{\centering\arraybackslash}X|>{\centering\arraybackslash}X|}
            \hline
            \textbf{Model} & \textbf{Benar/\allowbreak Total} & \textbf{Akurasi} & \textbf{Macro Precision} & \textbf{Macro Recall} & \textbf{Macro F1-score} \\
            \hline
            YOLOv8 Nano & 637/720 & 88,47\% & 89,03\% & 88,47\% & 88,63\% \\
            \hline
            YOLO11 Nano & 634/720 & 88,06\% & 88,30\% & 88,06\% & 88,12\% \\
            \hline
        \end{tabularx}
    \end{table}

    Berdasarkan Tabel \ref{tab:hasil-klasifikasi}, YOLOv8 Nano memperoleh akurasi sebesar 88,47\%, sedangkan YOLO11 Nano memperoleh akurasi sebesar 88,06\%. Selisih akurasi kedua model adalah 0,41\%, sehingga dapat dikatakan bahwa performa klasifikasi kedua model relatif sangat berdekatan. Pada skenario pengujian ini, YOLOv8 Nano memberikan hasil klasifikasi yang sedikit lebih tinggi dibandingkan YOLO11 Nano.

    Selain akurasi, nilai \textit{macro precision}, \textit{macro recall}, dan \textit{macro F1-score} juga menunjukkan bahwa YOLOv8 Nano sedikit lebih unggul. Nilai \textit{macro F1-score} YOLOv8 Nano sebesar 88,63\%, sedangkan YOLO11 Nano sebesar 88,12\%. Karena jumlah data uji pada setiap kelas seimbang, nilai akurasi dan nilai makro dapat digunakan bersama untuk melihat performa keseluruhan model secara lebih objektif.
\end{sectioncontent}

\section{Hasil Evaluasi per Kelas}
\begin{sectioncontent}
    \hspace{\parindent}Evaluasi per kelas dilakukan untuk mengetahui performa model pada masing-masing kategori penyakit daun tomat. Hasil evaluasi per kelas untuk YOLOv8 Nano ditampilkan pada Tabel \ref{tab:perkelas-yolov8}, sedangkan hasil evaluasi per kelas untuk YOLO11 Nano ditampilkan pada Tabel \ref{tab:perkelas-yolo11}.

    \begin{table}[H]
        \centering
        \caption{Hasil Evaluasi per Kelas pada YOLOv8 Nano}
        \label{tab:perkelas-yolov8}
        \begin{tabular}{|l|c|c|c|c|}
            \hline
            \textbf{Kelas} & \textbf{Precision} & \textbf{Recall} & \textbf{F1-score} & \textbf{Support} \\
            \hline
            Early Blight & 73,38\% & 85,00\% & 78,76\% & 120 \\
            \hline
            Healthy & 97,52\% & 98,33\% & 97,93\% & 120 \\
            \hline
            Yellow Curl Virus & 100,00\% & 90,00\% & 94,74\% & 120 \\
            \hline
            Mold Leaf & 86,96\% & 83,33\% & 85,11\% & 120 \\
            \hline
            Septoria Leaf Spot & 89,17\% & 89,17\% & 89,17\% & 120 \\
            \hline
            Late Blight & 87,18\% & 85,00\% & 86,08\% & 120 \\
            \hline
        \end{tabular}
    \end{table}

    \begin{table}[H]
        \centering
        \caption{Hasil Evaluasi per Kelas pada YOLO11 Nano}
        \label{tab:perkelas-yolo11}
        \begin{tabular}{|l|c|c|c|c|}
            \hline
            \textbf{Kelas} & \textbf{Precision} & \textbf{Recall} & \textbf{F1-score} & \textbf{Support} \\
            \hline
            Early Blight & 78,46\% & 85,00\% & 81,60\% & 120 \\
            \hline
            Healthy & 95,69\% & 92,50\% & 94,07\% & 120 \\
            \hline
            Yellow Curl Virus & 92,11\% & 87,50\% & 89,74\% & 120 \\
            \hline
            Mold Leaf & 89,29\% & 83,33\% & 86,21\% & 120 \\
            \hline
            Septoria Leaf Spot & 89,34\% & 90,83\% & 90,08\% & 120 \\
            \hline
            Late Blight & 84,92\% & 89,17\% & 86,99\% & 120 \\
            \hline
        \end{tabular}
    \end{table}

    Berdasarkan Tabel \ref{tab:perkelas-yolov8} dan Tabel \ref{tab:perkelas-yolo11}, kedua model menunjukkan performa yang cukup baik pada sebagian besar kelas. YOLOv8 Nano memperoleh nilai \textit{recall} tertinggi pada kelas Healthy sebesar 98,33\% dan nilai \textit{precision} tertinggi pada kelas Yellow Curl Virus sebesar 100,00\%. Sementara itu, YOLO11 Nano memperoleh performa yang baik pada kelas Healthy, Septoria Leaf Spot, dan Late Blight.

    Kelas Early Blight menjadi salah satu kelas yang relatif lebih sulit diklasifikasikan oleh kedua model, terutama pada YOLOv8 Nano yang memiliki nilai \textit{precision} lebih rendah dibandingkan kelas lain. Kondisi tersebut menunjukkan bahwa beberapa citra dari kelas lain masih diprediksi sebagai Early Blight. Hal ini dapat disebabkan oleh kemiripan pola visual antara Early Blight dengan beberapa penyakit daun lain, terutama pada tekstur bercak dan perubahan warna daun.
\end{sectioncontent}

\section{Analisis Confusion Matrix}
\begin{sectioncontent}
    \hspace{\parindent}\textit{Confusion matrix} digunakan untuk melihat distribusi prediksi benar dan salah pada masing-masing kelas. Untuk memudahkan penyajian tabel, digunakan singkatan kelas sebagai berikut: EB untuk Early Blight, HL untuk Healthy Leaf, YCV untuk Yellow Curl Virus, ML untuk Mold Leaf, SLS untuk Septoria Leaf Spot, dan LB untuk Late Blight.

    \begin{table}[H]
        \centering
        \caption{Confusion Matrix YOLOv8 Nano}
        \label{tab:cm-yolov8}
        \begin{tabular}{|l|c|c|c|c|c|c|}
            \hline
            \textbf{Aktual/Prediksi} & \textbf{EB} & \textbf{HL} & \textbf{YCV} & \textbf{ML} & \textbf{SLS} & \textbf{LB} \\
            \hline
            EB & 102 & 0 & 0 & 11 & 1 & 6 \\
            \hline
            HL & 0 & 118 & 0 & 1 & 0 & 1 \\
            \hline
            YCV & 12 & 0 & 108 & 0 & 0 & 0 \\
            \hline
            ML & 8 & 2 & 0 & 100 & 2 & 8 \\
            \hline
            SLS & 9 & 1 & 0 & 3 & 107 & 0 \\
            \hline
            LB & 8 & 0 & 0 & 0 & 10 & 102 \\
            \hline
        \end{tabular}
    \end{table}

    \begin{table}[H]
        \centering
        \caption{Confusion Matrix YOLO11 Nano}
        \label{tab:cm-yolo11}
        \begin{tabular}{|l|c|c|c|c|c|c|}
            \hline
            \textbf{Aktual/Prediksi} & \textbf{EB} & \textbf{HL} & \textbf{YCV} & \textbf{ML} & \textbf{SLS} & \textbf{LB} \\
            \hline
            EB & 102 & 0 & 2 & 7 & 3 & 6 \\
            \hline
            HL & 1 & 111 & 1 & 0 & 1 & 6 \\
            \hline
            YCV & 11 & 1 & 105 & 0 & 3 & 0 \\
            \hline
            ML & 5 & 2 & 5 & 100 & 2 & 6 \\
            \hline
            SLS & 2 & 2 & 1 & 5 & 109 & 1 \\
            \hline
            LB & 9 & 0 & 0 & 0 & 4 & 107 \\
            \hline
        \end{tabular}
    \end{table}

    \newpage

    Berdasarkan Tabel \ref{tab:cm-yolov8}, kesalahan prediksi terbesar pada YOLOv8 Nano terjadi pada kelas Yellow Curl Virus yang diprediksi sebagai Early Blight sebanyak 12 citra, kelas Early Blight yang diprediksi sebagai Mold Leaf sebanyak 11 citra, serta kelas Late Blight yang diprediksi sebagai Septoria Leaf Spot sebanyak 10 citra. Pada Tabel \ref{tab:cm-yolo11}, kesalahan prediksi terbesar terjadi pada kelas Yellow Curl Virus yang diprediksi sebagai Early Blight sebanyak 11 citra, serta kelas Late Blight yang diprediksi sebagai Early Blight sebanyak 9 citra.

    Kesalahan prediksi tersebut menunjukkan bahwa beberapa kelas penyakit daun tomat memiliki kemiripan pola visual. Penyakit seperti Early Blight, Septoria Leaf Spot, Mold Leaf, dan Late Blight sama-sama dapat menampilkan gejala bercak atau perubahan warna pada permukaan daun, sehingga model dapat mengalami kesulitan dalam membedakan kelas-kelas tersebut pada citra tertentu.
\end{sectioncontent}

\section{Hasil Pengujian Efisiensi pada Perangkat Edge}
\begin{sectioncontent}
    \hspace{\parindent}Pengujian efisiensi dilakukan pada Raspberry Pi 5 menggunakan model hasil export dalam format ONNX. Pengujian ini bertujuan untuk mengetahui waktu inferensi dan kemampuan pemrosesan model pada perangkat \textit{edge}. Ringkasan hasil pengujian efisiensi dapat dilihat pada Tabel \ref{tab:efisiensi-edge}.

    \begin{table}[H]
        \centering
        \caption{Perbandingan Efisiensi Inferensi pada Raspberry Pi 5}
        \label{tab:efisiensi-edge}
        \begin{tabular}{|l|c|c|c|c|c|}
            \hline
            \textbf{Model} & \textbf{Rata-rata} & \textbf{Median} & \textbf{Minimum} & \textbf{Maksimum} & \textbf{FPS} \\
            \hline
            YOLOv8 Nano & 7,78 ms & 7,53 ms & 7,43 ms & 22,10 ms & 128,52 \\
            \hline
            YOLO11 Nano & 7,86 ms & 7,60 ms & 7,48 ms & 20,19 ms & 127,18 \\
            \hline
        \end{tabular}
    \end{table}

    Berdasarkan Tabel \ref{tab:efisiensi-edge}, YOLOv8 Nano memiliki rata-rata waktu inferensi sebesar 7,78 ms, sedangkan YOLO11 Nano memiliki rata-rata waktu inferensi sebesar 7,86 ms. Dari sisi FPS, YOLOv8 Nano mencapai 128,52 FPS, sedangkan YOLO11 Nano mencapai 127,18 FPS. Perbedaan waktu inferensi antara kedua model relatif kecil, namun pada skenario pengujian ini YOLOv8 Nano sedikit lebih cepat dibandingkan YOLO11 Nano.

    Nilai maksimum waktu inferensi pada kedua model lebih tinggi dibandingkan nilai rata-rata dan median. Hal ini dapat terjadi karena adanya variasi beban sistem pada Raspberry Pi saat proses pengujian, seperti aktivitas sistem operasi, proses latar belakang, atau kondisi termal perangkat. Oleh karena itu, nilai rata-rata dan median digunakan secara bersamaan untuk memperoleh gambaran efisiensi inferensi yang lebih seimbang.
\end{sectioncontent}

\section{Pembahasan}
\begin{sectioncontent}
    \hspace{\parindent}Hasil pengujian menunjukkan bahwa YOLOv8 Nano memperoleh performa yang sedikit lebih baik dibandingkan YOLO11 Nano pada dataset dan skenario pengujian yang digunakan. Dari sisi akurasi, YOLOv8 Nano memperoleh nilai 88,47\%, sedangkan YOLO11 Nano memperoleh nilai 88,06\%. Selisih ini menunjukkan bahwa kedua model memiliki kemampuan klasifikasi yang sangat berdekatan, tetapi YOLOv8 Nano memberikan jumlah prediksi benar yang sedikit lebih banyak pada data uji.

    Dari sisi evaluasi per kelas, YOLOv8 Nano unggul pada kelas Healthy dan Yellow Curl Virus, terutama karena nilai \textit{recall} pada kelas Healthy dan nilai \textit{precision} pada kelas Yellow Curl Virus cukup tinggi. Sementara itu, YOLO11 Nano menunjukkan performa yang kompetitif pada kelas Early Blight, Septoria Leaf Spot, dan Late Blight. Perbedaan ini memperlihatkan bahwa masing-masing model memiliki karakteristik yang berbeda dalam mengenali pola visual pada citra daun tomat.

    Pada analisis \textit{confusion matrix}, beberapa kelas penyakit masih menunjukkan kesalahan klasifikasi yang cukup terlihat. Pada YOLOv8 Nano, kelas Yellow Curl Virus beberapa kali diprediksi sebagai Early Blight, sedangkan kelas Early Blight beberapa kali diprediksi sebagai Mold Leaf. Pada YOLO11 Nano, kesalahan prediksi juga terjadi pada kelas Yellow Curl Virus dan Late Blight yang diprediksi sebagai Early Blight. Kesalahan ini dapat terjadi karena beberapa penyakit daun tomat memiliki gejala visual yang mirip, seperti bercak pada daun, perubahan warna, dan pola kerusakan pada permukaan daun. Dengan demikian, kesalahan klasifikasi tidak hanya dipengaruhi oleh kemampuan model, tetapi juga oleh tingkat kemiripan visual antar kelas.

    Dari sisi efisiensi inferensi pada Raspberry Pi 5, YOLOv8 Nano juga sedikit lebih unggul dibandingkan YOLO11 Nano. Rata-rata waktu inferensi YOLOv8 Nano sebesar 7,78 ms, sedangkan YOLO11 Nano sebesar 7,86 ms. Perbedaan ini tidak terlalu besar, namun tetap menunjukkan bahwa pada lingkungan pengujian yang digunakan, YOLOv8 Nano mampu memberikan waktu inferensi yang sedikit lebih rendah.

    Berdasarkan kombinasi antara akurasi dan efisiensi, YOLOv8 Nano dapat dinilai sebagai model yang lebih optimal pada skenario penelitian ini. Namun, hasil tersebut tidak dapat digeneralisasi secara mutlak untuk seluruh dataset atau seluruh perangkat \textit{edge}. Kesimpulan yang dapat diambil adalah bahwa YOLOv8 Nano lebih unggul pada dataset, konfigurasi pelatihan, format model, dan perangkat pengujian yang digunakan dalam penelitian ini.
\end{sectioncontent}

\section{Keterbatasan Hasil Pengujian}
\begin{sectioncontent}
    \hspace{\parindent}Penelitian ini memiliki beberapa keterbatasan yang perlu diperhatikan dalam interpretasi hasil. Pertama, meskipun dataset pengujian yang digunakan sudah seimbang, penelitian ini masih menggunakan satu dataset utama sehingga hasilnya belum tentu sama apabila diterapkan pada dataset lain dengan kondisi pencahayaan, latar belakang, atau kualitas citra yang berbeda. Kedua, pengujian efisiensi dilakukan pada satu jenis perangkat \textit{edge}, yaitu Raspberry Pi 5, sehingga hasilnya belum tentu sama apabila diuji pada perangkat lain. Ketiga, pengujian dilakukan menggunakan model hasil export dalam format ONNX, sehingga performa dapat dipengaruhi oleh runtime yang digunakan. Keempat, hasil penelitian ini berfokus pada klasifikasi citra, bukan deteksi lokasi penyakit pada bagian daun.

    Dengan keterbatasan tersebut, hasil penelitian tetap dapat digunakan sebagai dasar perbandingan antara YOLOv8 Nano dan YOLO11 Nano dalam skenario klasifikasi penyakit daun tomat berbasis perangkat \textit{edge}. Namun, penelitian lanjutan dapat dilakukan dengan menambah variasi dataset, menguji perangkat \textit{edge} lain, serta membandingkan format model yang berbeda seperti NCNN atau TensorFlow Lite.
\end{sectioncontent}

\section{Ringkasan Hasil}
\begin{sectioncontent}
    \hspace{\parindent}Berdasarkan hasil pengujian yang telah dilakukan, dapat dirangkum bahwa YOLOv8 Nano memperoleh akurasi sebesar 88,47\% dengan rata-rata waktu inferensi 7,78 ms, sedangkan YOLO11 Nano memperoleh akurasi sebesar 88,06\% dengan rata-rata waktu inferensi 7,86 ms. Dari hasil tersebut, YOLOv8 Nano menunjukkan performa yang sedikit lebih baik dibandingkan YOLO11 Nano pada skenario penelitian ini.

    Secara umum, kedua model mampu menjalankan klasifikasi penyakit daun tomat dengan akurasi di atas 88\% dan waktu inferensi di bawah 10 ms pada Raspberry Pi 5. Hal ini menunjukkan bahwa kedua model layak dipertimbangkan untuk implementasi klasifikasi citra pada perangkat \textit{edge}. Namun, berdasarkan hasil pengujian yang diperoleh, YOLOv8 Nano menjadi model yang lebih direkomendasikan dalam penelitian ini karena memberikan keseimbangan yang sedikit lebih baik antara akurasi dan efisiensi inferensi.
\end{sectioncontent}
